\begin{proof}[Proof of \cref{obj:thm:sharp-honest-frontier}]
\begin{enumerate}
\item Fix the joint calibration constants \(\varepsilon,M\) supplied by
\cref{obj:lem:calibrated-support}, and use the corresponding positive constants
\(\kappa,C_{\mathrm{mix}}\) from \cref{obj:lem:original-record-mixtures}. Define
\[
L_{\mathrm{mix}}
=
\left\lceil
\max\{1,\kappa^{-1},100C_{\mathrm{mix}}\varepsilon^4\}
\right\rceil,
\qquad
d_{\mathrm{ceil}}=(L_{\mathrm{mix}}+1)^{-1/2}.
\]
These are finite absolute constants, and the ceiling gives
\[
L_{\mathrm{mix}}\ge1,\qquad
\frac1{L_{\mathrm{mix}}}\le\kappa,\qquad
\frac{C_{\mathrm{mix}}\varepsilon^4}{L_{\mathrm{mix}}}
\le\frac1{100},\qquad
d_{\mathrm{ceil}}>0.
\]
For the argument below, write the worst expected length of an admissible procedure as
\[
\mathcal W_n(\alpha,\beta;r,I)
=
\sup_{P\in\mathcal M_r(\alpha,\beta)}
\mathbb E_{P^{\otimes n}\otimes\lambda}
|I(\mathcal O,U)|.
\]
We work throughout with \((\alpha,\beta)\in\mathcal D\), \(n\ge1\),
\(r\in[0,1/2]\), and \(I\in\mathcal A_n(\alpha,\beta)\), unless a narrower
range is specified.
% lean: CausalSmith.Stat.LogoddsLowsmoothFrontier.frontier_mixture_lower_rates

\item We first establish the rank estimates used in both mixture constructions.
For every real \(\gamma_{\mathrm{rank}}\in(0,1/2]\), define
\[
k_{\mathrm{low}}(\gamma_{\mathrm{rank}},n)
=
\left\lceil
L_{\mathrm{mix}}n^{2/(2\gamma_{\mathrm{rank}}+1)}
\right\rceil.
\]
Since \(n^{2/(2\gamma_{\mathrm{rank}}+1)}\ge1\), the ceiling inequalities yield
\[
1\le k_{\mathrm{low}}(\gamma_{\mathrm{rank}},n),\qquad
L_{\mathrm{mix}}n^{2/(2\gamma_{\mathrm{rank}}+1)}
\le k_{\mathrm{low}}(\gamma_{\mathrm{rank}},n)
\le
(L_{\mathrm{mix}}+1)n^{2/(2\gamma_{\mathrm{rank}}+1)}.
\]
The exponent \(2/(2\gamma_{\mathrm{rank}}+1)\) is at least one, so
\[
L_{\mathrm{mix}}n
\le
L_{\mathrm{mix}}n^{2/(2\gamma_{\mathrm{rank}}+1)}
\le k_{\mathrm{low}}(\gamma_{\mathrm{rank}},n),
\qquad
\frac{n}{k_{\mathrm{low}}(\gamma_{\mathrm{rank}},n)}
\le\frac1{L_{\mathrm{mix}}}.
\]
Moreover, \(L_{\mathrm{mix}}^{2\gamma_{\mathrm{rank}}+1}\ge L_{\mathrm{mix}}\).
Raising the lower rank bound to the positive power \(2\gamma_{\mathrm{rank}}+1\) therefore gives
\[
k_{\mathrm{low}}(\gamma_{\mathrm{rank}},n)^{2\gamma_{\mathrm{rank}}+1}
\ge
L_{\mathrm{mix}}^{2\gamma_{\mathrm{rank}}+1}n^2
\ge L_{\mathrm{mix}}n^2,
\]
and hence
\[
n^2k_{\mathrm{low}}(\gamma_{\mathrm{rank}},n)^{-(2\gamma_{\mathrm{rank}}+1)}
\le\frac1{L_{\mathrm{mix}}}.
\]
Finally, raising the upper rank bound to the negative power
\(-\gamma_{\mathrm{rank}}\), and using
\(\gamma_{\mathrm{rank}}\le1/2\), gives
\[
\begin{aligned}
k_{\mathrm{low}}(\gamma_{\mathrm{rank}},n)^{-\gamma_{\mathrm{rank}}}
&\ge
(L_{\mathrm{mix}}+1)^{-\gamma_{\mathrm{rank}}}
n^{-2\gamma_{\mathrm{rank}}/(2\gamma_{\mathrm{rank}}+1)}\\
&\ge
d_{\mathrm{ceil}}\,
n^{-2\gamma_{\mathrm{rank}}/(2\gamma_{\mathrm{rank}}+1)}.
\end{aligned}
\]
All these estimates include \(n=1\).
% lean: CausalSmith.Stat.LogoddsLowsmoothFrontier.lower_rank_ceiling_bounds
% lean: CausalSmith.Stat.LogoddsLowsmoothFrontier.lower_rank_budget
% lean: CausalSmith.Stat.LogoddsLowsmoothFrontier.lower_rank_separation

\item Suppose first that \(\alpha+\beta\le1/2\). Define the mixed rank and amplitudes by
\[
k_{\mathrm{mixed}}=k_{\mathrm{low}}(\alpha+\beta,n),
\qquad
\eta=\varepsilon k_{\mathrm{mixed}}^{-\alpha},
\qquad
\zeta=\varepsilon k_{\mathrm{mixed}}^{-\beta}.
\]
The domain conditions imply \(0<\alpha+\beta\le1/2\) and
\(\alpha,\beta>0\). Since \(k_{\mathrm{mixed}}\ge1\),
\[
0<\eta\le\varepsilon,\qquad 0<\zeta\le\varepsilon,
\qquad
\eta\zeta=\varepsilon^2k_{\mathrm{mixed}}^{-(\alpha+\beta)}.
\]
The rank estimates and the calibration of \(L_{\mathrm{mix}}\) give
\[
\frac{n}{k_{\mathrm{mixed}}}\le\kappa
\]
and, by the laws of real powers,
\[
\begin{aligned}
C_{\mathrm{mix}}\frac{n^2}{k_{\mathrm{mixed}}}\eta^2\zeta^2
&=
C_{\mathrm{mix}}\varepsilon^4
n^2k_{\mathrm{mixed}}^{-(2\alpha+2\beta+1)}\\
&\le
\frac{C_{\mathrm{mix}}\varepsilon^4}{L_{\mathrm{mix}}}
\le\frac1{100}.
\end{aligned}
\]
Thus \cref{obj:lem:original-record-mixtures}, applied to the mixed priors of
\cref{obj:def:continuous-calibrated-priors}, supplies
\[
\mathcal H^2(Q_0,Q_1)\le\frac1{100}.
\]

We give the testing-to-length argument needed here and for the fair priors.
Let \(Q_{\mathrm{eval}}\) and \(Q_{\mathrm{alt}}\) be uniform finite mixtures
of original-record sample laws. More precisely, for
\(\iota_{\mathrm{prior}}\in\{\mathrm{eval},\mathrm{alt}\}\), take an integer \(k_{\iota_{\mathrm{prior}}}\ge0\), define
\[
\Sigma_{\iota_{\mathrm{prior}}}=\{-1,1\}^{k_{\iota_{\mathrm{prior}}}+1},
\qquad
Q_{\iota_{\mathrm{prior}}}=\frac1{|\Sigma_{\iota_{\mathrm{prior}}}|}
\sum_{\sigma\in\Sigma_{\iota_{\mathrm{prior}}}}P_{\iota_{\mathrm{prior}},\sigma}^{\otimes n},
\]
and suppose that their constituents have common effects
\[
\theta(P_{\mathrm{eval},\sigma})=t_{\mathrm{eval}},
\qquad
\theta(P_{\mathrm{alt},\sigma})=t_{\mathrm{alt}}.
\]
Assume that every evaluation constituent belongs to
\(\mathcal M_r(\alpha,\beta)\), every alternative constituent belongs to
\(\mathcal M(\alpha,\beta)\), and all constituents have uniform covariate law.

Let \(\tau_{\mathrm{lab}}\) be the uniform probability law on the four binary
label pairs, and define the original-record reference and mixture densities by
\[
\tau_{\mathrm{lab}}
=\frac14\sum_{(a_{\mathrm{lab}},y_{\mathrm{lab}})\in\{0,1\}^2}
\delta_{(a_{\mathrm{lab}},y_{\mathrm{lab}})},
\qquad
R_{\mathrm{sample}}=(\lambda\otimes\tau_{\mathrm{lab}})^{\otimes n},
\qquad
q_{\iota_{\mathrm{prior}}}=\frac{dQ_{\iota_{\mathrm{prior}}}}{dR_{\mathrm{sample}}}.
\]
These densities are nonnegative and integrate to one. With total variation defined by
\[
\operatorname{TV}(Q_{\mathrm{eval}},Q_{\mathrm{alt}})
=
\frac12\int
|q_{\mathrm{eval}}-q_{\mathrm{alt}}|\,dR_{\mathrm{sample}},
\]
the squared Hellinger distance satisfies
\[
\mathcal H^2(Q_{\mathrm{eval}},Q_{\mathrm{alt}})
=
\int
(\sqrt{q_{\mathrm{eval}}}-\sqrt{q_{\mathrm{alt}}})^2\,dR_{\mathrm{sample}}
=
2-2\int\sqrt{q_{\mathrm{eval}}q_{\mathrm{alt}}}\,dR_{\mathrm{sample}}.
\]
Cauchy--Schwarz gives
\[
\int\sqrt{q_{\mathrm{eval}}q_{\mathrm{alt}}}\,dR_{\mathrm{sample}}\le1,
\qquad
\int(\sqrt{q_{\mathrm{eval}}}+\sqrt{q_{\mathrm{alt}}})^2
\,dR_{\mathrm{sample}}\le4.
\]
Factoring the density difference and applying Cauchy--Schwarz again yields
\[
\operatorname{TV}(Q_{\mathrm{eval}},Q_{\mathrm{alt}})
\le
\frac12
\sqrt{\mathcal H^2(Q_{\mathrm{eval}},Q_{\mathrm{alt}})}
\left[
\int(\sqrt{q_{\mathrm{eval}}}+\sqrt{q_{\mathrm{alt}}})^2
\,dR_{\mathrm{sample}}
\right]^{1/2}
\le
\sqrt{\mathcal H^2(Q_{\mathrm{eval}},Q_{\mathrm{alt}})}.
\]
The product densities with respect to
\(R_{\mathrm{sample}}\otimes\lambda\) are the same densities, independent of
the seed. Integrating over that seed proves
\[
\operatorname{TV}(Q_{\mathrm{eval}}\otimes\lambda,
Q_{\mathrm{alt}}\otimes\lambda)
=
\operatorname{TV}(Q_{\mathrm{eval}},Q_{\mathrm{alt}}).
\]
Consequently, a squared Hellinger bound of \(1/100\) gives total variation
at most \(1/10\) for the sample-and-seed experiments.

Define the measurable coverage events
\[
\mathcal E_{\mathrm{eval}}
=\{(\mathcal O,U):t_{\mathrm{eval}}\in I(\mathcal O,U)\},
\qquad
\mathcal E_{\mathrm{alt}}
=\{(\mathcal O,U):t_{\mathrm{alt}}\in I(\mathcal O,U)\}.
\]
Global honesty and finite averaging imply
\[
(Q_{\mathrm{eval}}\otimes\lambda)(\mathcal E_{\mathrm{eval}})\ge\frac9{10},
\qquad
(Q_{\mathrm{alt}}\otimes\lambda)(\mathcal E_{\mathrm{alt}})\ge\frac9{10}.
\]
Transferring the second event by total variation, then using that a union has
probability at most one, gives
\[
(Q_{\mathrm{eval}}\otimes\lambda)
(\mathcal E_{\mathrm{eval}}\cap\mathcal E_{\mathrm{alt}})
\ge\frac9{10}+\left(\frac9{10}-\frac1{10}\right)-1
=\frac7{10}.
\]
Connectedness of the returned interval gives the pointwise inequality
\[
|I(\mathcal O,U)|
\ge
|t_{\mathrm{alt}}-t_{\mathrm{eval}}|\,
\mathbf1_{\mathcal E_{\mathrm{eval}}\cap\mathcal E_{\mathrm{alt}}}
(\mathcal O,U).
\]
Integration and finite averaging therefore show
\[
\begin{aligned}
\frac7{10}|t_{\mathrm{alt}}-t_{\mathrm{eval}}|
&\le
\mathbb E_{Q_{\mathrm{eval}}\otimes\lambda}|I(\mathcal O,U)|\\
&=
\frac1{|\Sigma_{\mathrm{eval}}|}
\sum_{\sigma\in\Sigma_{\mathrm{eval}}}
\mathbb E_{P_{\mathrm{eval},\sigma}^{\otimes n}\otimes\lambda}
|I(\mathcal O,U)|\\
&\le\mathcal W_n(\alpha,\beta;r,I).
\end{aligned}
\]
The procedure's measurability and bounded interval length justify these
integrals. This argument includes independently randomized procedures.
% lean: CausalSmith.Stat.LogoddsLowsmoothFrontier.finitePrior_hellinger_length_lower

For the mixed priors, \cref{obj:lem:calibrated-support} supplies model
membership and the effects
\[
\theta(P_{0,\sigma})=0,\qquad
\theta(P_{1,\sigma})=32\eta\zeta.
\]
The null constituents belong to every evaluation slice because \(r\ge0\).
Applying the preceding reduction with \(Q_0\) as evaluation prior gives
\[
\mathcal W_n(\alpha,\beta;r,I)
\ge\frac7{10}\,32\eta\zeta
\ge
\frac7{10}\,32\varepsilon^2d_{\mathrm{ceil}}\,
n^{-2(\alpha+\beta)/(2\alpha+2\beta+1)}.
\]
Since the denominator is positive,
\[
\frac{2(\alpha+\beta)}{2\alpha+2\beta+1}\le\frac12
\quad\Longleftrightarrow\quad
\alpha+\beta\le\frac12.
\]
Thus, defining
\[
c_{\mathrm{mix}}=\frac7{10}\,32\varepsilon^2d_{\mathrm{ceil}}>0,
\]
we have proved
\[
c_{\mathrm{mix}}n^{-\mathsf a(\alpha,\beta)}
\le\mathcal W_n(\alpha,\beta;r,I)
\qquad(\alpha+\beta\le1/2).
\]
% lean: CausalSmith.Stat.LogoddsLowsmoothFrontier.mixed_prior_length_lower
% lean: CausalSmith.Stat.LogoddsLowsmoothFrontier.frontier_mixture_lower_rates

\item For the radius-dependent lower bound, suppose \(r>0\) and define
\[
k_{\mathrm{fair}}=k_{\mathrm{low}}(2\beta,n),
\qquad
\delta=\varepsilon k_{\mathrm{fair}}^{-\beta},
\qquad
t_{\mathrm{fair}}=\frac r2.
\]
Here \(0<2\beta<1/2\), \(0<\delta\le\varepsilon\le1/100\), and
\(t_{\mathrm{fair}}\in(0,1/4]\). The rank estimates imply
\[
\frac{n}{k_{\mathrm{fair}}}\le\kappa,
\qquad
C_{\mathrm{mix}}\frac{n^2}{k_{\mathrm{fair}}}\delta^4
=
C_{\mathrm{mix}}\varepsilon^4
n^2k_{\mathrm{fair}}^{-(4\beta+1)}
\le\frac1{100}.
\]
By \cref{obj:lem:original-record-mixtures},
\[
\mathcal H^2\!\left(
Q_{t_{\mathrm{fair}}},
P_{*,t_{\mathrm{fair}},\delta}^{\otimes n}
\right)\le\frac1{100}.
\]
The fair support assertion in \cref{obj:lem:calibrated-support} supplies
model membership and
\[
\theta(P_{t_{\mathrm{fair}},\sigma})=t_{\mathrm{fair}},
\qquad
\theta(P_{*,t_{\mathrm{fair}},\delta})
=\mathsf T(t_{\mathrm{fair}},\delta).
\]
The random constituents belong to the evaluation slice because
\(0<t_{\mathrm{fair}}\le r\). The comparator can be regarded as a constant
finite prior: averaging identical copies leaves its sample law unchanged.
By \cref{obj:lem:exact-calibrations},
\[
\frac{t_{\mathrm{fair}}}{2}
\le\mathsf T(t_{\mathrm{fair}},\delta)\le t_{\mathrm{fair}},
\qquad
t_{\mathrm{fair}}-\mathsf T(t_{\mathrm{fair}},\delta)
\ge3t_{\mathrm{fair}}\delta^2.
\]
Using the random fair prior as evaluation prior in the testing reduction,
and then the rank separation estimate at \(\gamma_{\mathrm{rank}}=2\beta\),
we obtain
\[
\begin{aligned}
\mathcal W_n(\alpha,\beta;r,I)
&\ge
\frac7{10}
\bigl[t_{\mathrm{fair}}-\mathsf T(t_{\mathrm{fair}},\delta)\bigr]\\
&\ge
\frac7{10}\,\frac32 r\varepsilon^2k_{\mathrm{fair}}^{-2\beta}\\
&\ge
\frac7{10}\,\frac32\varepsilon^2d_{\mathrm{ceil}}\,
r n^{-4\beta/(4\beta+1)}.
\end{aligned}
\]
Define
\[
c_{\mathrm{radius}}
=\frac7{10}\,\frac32\varepsilon^2d_{\mathrm{ceil}}>0.
\]
We have therefore proved
\[
c_{\mathrm{radius}}r n^{-\mathsf b(\beta)}
\le\mathcal W_n(\alpha,\beta;r,I)
\qquad(0<r\le1/2).
\]
% lean: CausalSmith.Stat.LogoddsLowsmoothFrontier.fair_prior_length_lower
% lean: CausalSmith.Stat.LogoddsLowsmoothFrontier.frontier_mixture_lower_rates

\item Choose the positive finite absolute upper constant supplied by
\cref{obj:thm:finite-honest-upper}, and set
\[
C=K_0.
\]
For every admissible engine and naming policy, that result supplies
\[
I^{\mathrm{up}}_{n,\mathsf E,\mathsf N}(\alpha,\beta)
\in\mathcal A_n(\alpha,\beta),
\]
together with
\[
\mathbb E_{P^{\otimes n}\otimes\lambda}
\left|I^{\mathrm{up}}_{n,\mathsf E,\mathsf N}(\alpha,\beta)(\mathcal O,U)\right|
\le C\rho_n(\alpha,\beta;r)
\qquad(P\in\mathcal M_r(\alpha,\beta)).
\]
Here the identification of the diagnostic envelope with \(\rho_n\) is
\cref{obj:def:frontier-handle}. Taking the supremum over the slice and using
the infimum definition in \cref{obj:def:length-objective} gives
\[
J_n(\alpha,\beta;r)
\le
\mathcal W_n\!\left(
\alpha,\beta;r,I^{\mathrm{up}}_{n,\mathsf E,\mathsf N}(\alpha,\beta)
\right)
\le C\rho_n(\alpha,\beta;r).
\]
The slice is nonempty: the law
\[
P_{\mathrm{null}}=\lambda\otimes\tau_{\mathrm{lab}}
\]
has
\[
e(P_{\mathrm{null}},x)=\mu_0(P_{\mathrm{null}},x)
=\mu_1(P_{\mathrm{null}},x)=\frac12,
\qquad
g(P_{\mathrm{null}},x)=\nu(P_{\mathrm{null}},x)
=\theta(P_{\mathrm{null}})=0,
\]
and hence satisfies the model conditions and belongs to every slice with
\(r\ge0\).

The upper result also supplies Borel measurability, global \(90\%\) coverage,
and the ordered rational closed output in the effect envelope. The
construction in \cref{obj:def:upper-handle} depends on the fixed engine,
public inputs, selected names, and original records, and therefore satisfies
\[
I^{\mathrm{up}}_{n,\mathsf E,\mathsf N}(\alpha,\beta)(\mathcal O,U)
=
I^{\mathrm{up}}_{n,\mathsf E,\mathsf N}(\alpha,\beta)(\mathcal O).
\]
Its input list makes it radius independent. For \(n\ge2\),
\cref{obj:lem:fixed-budget-realization} supplies successful termination
within the two prescribed budgets and finite rational evaluation. For
\(n=1\), the construction returns \([-1/2,1/2]\) directly.
% lean: CausalSmith.Stat.LogoddsLowsmoothFrontier.frontier_upper_comparisons

\item Define the common lower constants by
\[
c_{\mathrm{numerator}}
=\min\{c_{\mathrm{mix}},3/16\},
\qquad
c=\frac12\min\{c_{\mathrm{numerator}},c_{\mathrm{radius}}\}.
\]
Both constants are strictly positive. If \(\alpha+\beta\le1/2\), the mixed
bound gives
\[
c_{\mathrm{numerator}}n^{-\mathsf a(\alpha,\beta)}
\le
c_{\mathrm{mix}}n^{-\mathsf a(\alpha,\beta)}
\le\mathcal W_n(\alpha,\beta;r,I).
\]
If \(\alpha+\beta>1/2\), positivity of \(2\alpha+2\beta+1\) gives
\[
\frac12\le\frac{2(\alpha+\beta)}{2\alpha+2\beta+1},
\qquad
\mathsf a(\alpha,\beta)=\frac12.
\]
By \cref{obj:lem:parametric-null-floor} and the infimum definition of \(J_n\),
\[
c_{\mathrm{numerator}}n^{-\mathsf a(\alpha,\beta)}
\le
\frac3{16}n^{-1/2}
\le J_n(\alpha,\beta;r)
\le\mathcal W_n(\alpha,\beta;r,I).
\]
This proves the numerator bound throughout \(\mathcal D\), with equality
\(\alpha+\beta=1/2\) included in the first case.
% lean: CausalSmith.Stat.LogoddsLowsmoothFrontier.frontier_numerator_lower_of_rough

For the radius bound, \(r>0\) is covered above. At \(r=0\), nonnegativity of
interval length gives
\[
c_{\mathrm{radius}}\,0\,n^{-\mathsf b(\beta)}
=0\le\mathcal W_n(\alpha,\beta;0,I).
\]
Thus
\[
c_{\mathrm{radius}}r n^{-\mathsf b(\beta)}
\le\mathcal W_n(\alpha,\beta;r,I)
\qquad(0\le r\le1/2).
\]
% lean: CausalSmith.Stat.LogoddsLowsmoothFrontier.frontier_radius_lower_of_positive

Each summand of \(\rho_n\) is nonnegative. Reducing both coefficients to
their minimum and adding the two lower inequalities yields
\[
\begin{aligned}
c\rho_n(\alpha,\beta;r)
&=
\frac{\min\{c_{\mathrm{numerator}},c_{\mathrm{radius}}\}}2
\left[n^{-\mathsf a(\alpha,\beta)}
+r n^{-\mathsf b(\beta)}\right]\\
&\le
\frac12\left[
\mathcal W_n(\alpha,\beta;r,I)
+\mathcal W_n(\alpha,\beta;r,I)
\right]\\
&=\mathcal W_n(\alpha,\beta;r,I).
\end{aligned}
\]
This is the asserted converse for every honest procedure, including
independently randomized procedures.
% lean: CausalSmith.Stat.LogoddsLowsmoothFrontier.frontier_rate_lower_of_two

\item The honest class is nonempty: the constant procedure returning
\([-1/2,1/2]\) is Borel and covers every ambient-model effect by
\cref{obj:ass:effect-envelope}. Its length is one, while all procedure
objectives are nonnegative. Taking the infimum of the preceding
all-procedure lower bound is therefore legitimate and gives
\[
c\rho_n(\alpha,\beta;r)\le J_n(\alpha,\beta;r).
\]
Combining this with the upper comparison proves
\[
c\rho_n(\alpha,\beta;r)
\le J_n(\alpha,\beta;r)
\le
\mathcal W_n\!\left(
\alpha,\beta;r,I^{\mathrm{up}}_{n,\mathsf E,\mathsf N}(\alpha,\beta)
\right)
\le C\rho_n(\alpha,\beta;r).
\]
For the upper procedure, seed independence implies, for every law \(P\),
\[
\mathbb E_{P^{\otimes n}\otimes\lambda}
\left|I^{\mathrm{up}}_{n,\mathsf E,\mathsf N}(\alpha,\beta)(\mathcal O,U)\right|
=
\mathbb E_{P^{\otimes n}}
\left|I^{\mathrm{up}}_{n,\mathsf E,\mathsf N}(\alpha,\beta)(\mathcal O)\right|,
\]
which gives the displayed comparison in the statement. The constants were
chosen independently of the engine and naming policy.

By \cref{obj:lem:concrete-arithmetic-engine},
\(\mathsf E_0\) is admissible, and
\cref{obj:lem:fixed-budget-realization} supplies admissibility of
\(\mathsf N_0\). Specializing the universal comparison and operational
properties to this pair gives
\[
I_n^\star=I^{\mathrm{up}}_{n,\mathsf E_0,\mathsf N_0},
\qquad c_\star=c,\qquad C_\star=C,\qquad n_0=1.
\]
% lean: CausalSmith.Stat.LogoddsLowsmoothFrontier.universalFrontier_of_procedure_lower

\item Let \(\mathcal K\subseteq\mathcal D\) be nonempty and compact. The constant
functions with values \(c\) and \(1\) have singleton images on \(\mathcal K\), so
\[
\inf_{(\alpha,\beta)\in\mathcal K}c=c>0,
\qquad
\sup_{(\alpha,\beta)\in\mathcal K}1=1.
\]
The constant upper bound gives
\[
\sup_{(\alpha,\beta)\in\mathcal K}C\le C<\infty.
\]
These are the required compact-uniform conditions for the same constants
and threshold.
% lean: CausalSmith.Stat.LogoddsLowsmoothFrontier.frontier_compact_uniform

\item It remains to establish the exponent identities. The domain conditions give
\[
\alpha>0,\qquad
2\alpha+2\beta+1>0,\qquad
4\beta+1>0.
\]
As already verified by comparison of the two entries of the minimum,
\[
\mathsf a(\alpha,\beta)
=
\frac{2(\alpha+\beta)}{2\alpha+2\beta+1}
\quad\text{if }\alpha+\beta\le\frac12,
\qquad
\mathsf a(\alpha,\beta)=\frac12
\quad\text{if }\alpha+\beta\ge\frac12.
\]
In the first range, subtraction over the positive common denominator gives
\[
\begin{aligned}
\mathsf a(\alpha,\beta)-\mathsf b(\beta)
&=
\frac{
2(\alpha+\beta)(4\beta+1)
-4\beta(2\alpha+2\beta+1)}
{(2\alpha+2\beta+1)(4\beta+1)}\\
&=
\frac{2(\alpha-\beta)}
{(2\alpha+2\beta+1)(4\beta+1)}.
\end{aligned}
\]
In the second range,
\[
\mathsf a(\alpha,\beta)-\mathsf b(\beta)
=
\frac12-\frac{4\beta}{4\beta+1}
=
\frac{1-4\beta}{2(4\beta+1)}.
\]
If \(\alpha+\beta\le1/2\), positivity follows from \(\alpha>\beta\) and
the positive denominators. If \(\alpha+\beta>1/2\), it follows from
\(\beta<1/4\) and the positive denominator. Thus
\[
\mathsf a(\alpha,\beta)-\mathsf b(\beta)>0
\qquad((\alpha,\beta)\in\mathcal D).
\]
At \(\alpha+\beta=1/2\), both branch identities apply to the same difference,
and consequently
\[
\frac{2(\alpha-\beta)}
{(2\alpha+2\beta+1)(4\beta+1)}
=
\mathsf a(\alpha,\beta)-\mathsf b(\beta)
=
\frac{1-4\beta}{2(4\beta+1)}.
\]
This establishes agreement at the boundary and completes the proof.
% lean: CausalSmith.Stat.LogoddsLowsmoothFrontier.frontier_elbow_identities
\end{enumerate}
\end{proof}
